\chapterimage{figs/chapterhead.png}
\chapter{Expression Language Reference}
\label{chap:expression-language-reference}

This chapter documents the expression language as it exists in the current
repository. The scope is intentionally strict: only constructs that are
implemented in the scanner, grammar, driver, and unit tests are treated as
part of the reference. Anything else is listed as a limit.

\noindent\textbf{Objective.} Provide the reference description of the current
simulator expression language.

\noindent\textbf{Audience.} Developers and advanced users who need the
supported syntax surface, the current semantic model, and the known gaps.

\noindent\textbf{Scope.} Cover lexical forms, numbers, identifiers,
operators, precedence, functions, distributions, model references, indexed
access, errors, examples, and unsupported syntax.

\noindent\textbf{Status.} Confirmed by code inspection and unit tests. The
language is model-aware, sparse-indexed, and deliberately conservative about
unsupported constructs.

The chapter follows Figure~\ref{fig:expression-precedence-map},
Figure~\ref{fig:expression-symbol-resolution},
Figure~\ref{fig:expression-context-model}, and
Figure~\ref{fig:expression-error-anatomy}.

\section{How the language is resolved}
\label{sec:expression-reference-resolution}

The expression language is not a free-form string evaluator. The scanner first
tries to resolve a token against the live model. Only if that lookup succeeds
does the parser treat the text as a model reference. If it does not succeed,
the scanner falls back to numbers, operators, keywords, or an illegal-token
result.

That means the language is contextual:
\begin{itemize}
\item the same symbol can mean different things in different models;
\item the current event and current entity can influence the value of an
      expression;
\item model-backed names are resolved before the parser sees them as plain
      text.
\end{itemize}

\begin{figure}[htbp]
\centering
% FIGURE-SPEC-BEGIN
% ID: fig:expression-precedence-map
% Source of truth:
%   - source/parser/parserBisonFlex/bisonparser.yy
%   - source/tests/unit/test_parser_expressions.cpp
% Purpose:
%   Show the precedence layers used by the current expression grammar.
% Required visual content:
%   - literals and parentheses
%   - unary plus and minus
%   - power as a right-associative layer
%   - multiplicative, additive, relational, and logical layers
%   - command forms shown as side branches, not arithmetic operators
% Accessibility description:
%   A layered precedence chart that explains how the parser groups expressions before evaluation.
% Validation:
%   The chart must match the grammar nonterminals and the tested precedence behavior.
% FIGURE-SPEC-END
\figureplaceholder{figs/generated/expression-precedence-map}
\caption{Expression precedence and grouping in the current grammar.}
\label{fig:expression-precedence-map}
\end{figure}

\section{Lexical surface}
\label{sec:expression-reference-lexical-surface}

The scanner recognizes a small, explicit surface:
\begin{itemize}
\item decimal numbers such as \texttt{12} or \texttt{12.5};
\item hexadecimal numbers such as \texttt{0x10};
\item booleans \texttt{true} and \texttt{false};
\item arithmetic and relational symbols;
\item bracketed index syntax;
\item keywords for commands, functions, and model-backed symbols;
\item model names that are resolved at scan time.
\end{itemize}

Identifiers are not accepted as arbitrary text. The scanner checks the live
model for a matching data definition or simulation symbol. When it finds one,
it emits the corresponding token for the parser. When it does not, the token
becomes illegal and the parser reports an error.

There are no quoted string literals in the current scanner. The language is
expression-first, not string-literal-first. The only string-oriented keywords
currently declared by the scanner are \texttt{val}, \texttt{eval}, and
\texttt{leng}, and those tokens are not wired into active grammar productions
yet.

\section{Numbers and booleans}
\label{sec:expression-reference-numbers-booleans}

Numbers are accepted as decimal or hexadecimal literals. Decimal parsing
includes floating-point notation. Boolean literals are mapped to numeric
values:
\begin{itemize}
\item \texttt{true} becomes \texttt{1};
\item \texttt{false} becomes \texttt{0}.
\end{itemize}

This numeric treatment is consistent with the rest of the grammar, where
logical forms ultimately evaluate to numeric values too.

\section{Operators and precedence}
\label{sec:expression-reference-operators}

The current precedence structure is:
\begin{itemize}
\item parentheses;
\item unary \texttt{+} and \texttt{-};
\item power \texttt{\^{}} with right associativity;
\item multiplication and division;
\item addition and subtraction;
\item relational operators \texttt{<}, \texttt{>}, \texttt{<=},
      \texttt{>=}, \texttt{==}, \texttt{!=}, and \texttt{<>};
\item logical \texttt{not}, \texttt{and}, \texttt{nand}, \texttt{xor}, and
      \texttt{or}.
\end{itemize}

The tests confirm that precedence is real and not merely documented:
\texttt{1+2*3} evaluates to \texttt{7}, \texttt{(1+2)*3} evaluates to
\texttt{9}, and \verb|2^3^2| evaluates to \texttt{512}, which proves that
power is right-associative in the current grammar.

\section{Functions}
\label{sec:expression-reference-functions}

The active grammar currently recognizes several function families.

\subsection{Math and trigonometric functions}
\label{subsec:expression-reference-math-trig}

The following functions are implemented in the grammar:
\begin{itemize}
\item \texttt{sin}, \texttt{cos};
\item \texttt{round}, \texttt{trunc}, \texttt{frac};
\item \texttt{exp}, \texttt{sqrt}, \texttt{log}, \texttt{ln};
\item \texttt{mod}, \texttt{min}, \texttt{max}.
\end{itemize}

The current unit tests directly reproduce several of these:
\texttt{round(3.6)} returns \texttt{4}, \texttt{trunc(3.9)} returns
\texttt{3}, \texttt{frac(3.9)} returns \texttt{0.9}, and \texttt{sqrt(9)}
returns \texttt{3}.

\subsection{Probability distributions}
\label{subsec:expression-reference-probability}

The current stochastic functions are:
\begin{itemize}
\item \texttt{rnd};
\item \texttt{expo};
\item \texttt{norm};
\item \texttt{unif};
\item \texttt{weib};
\item \texttt{logn};
\item \texttt{gamm};
\item \texttt{erla};
\item \texttt{tria};
\item \texttt{beta};
\item \texttt{disc}.
\end{itemize}

The implemented distributions call through \texttt{Sampler\_if}. The unit
tests reproduce valid and invalid cases for \texttt{unif}, \texttt{tria},
\texttt{expo}, \texttt{weib}, \texttt{norm}, \texttt{logn},
\texttt{gamm}, \texttt{erla}, and \texttt{beta}. The invalid cases are
important because they prove that the parser reports recoverable errors rather
than silently accepting bad parameters.

\subsection{Kernel and model-aware functions}
\label{subsec:expression-reference-kernel}

The current grammar also exposes simulation and model-dependent helpers:
\begin{itemize}
\item \texttt{tnow}, \texttt{tfin}, \texttt{maxrep}, \texttt{numrep};
\item \texttt{ident}, \texttt{entitieswip};
\item \texttt{tavg}, \texttt{count}.
\end{itemize}

These values come from the live simulation or from model data definitions.
They are not pure arithmetic functions.

\subsection{Queue, resource, set, variable, and formula references}
\label{subsec:expression-reference-model-symbols}

The grammar includes model-backed references for:
\begin{itemize}
\item attributes;
\item variables;
\item formulas;
\item queues;
\item resources;
\item sets;
\item counters;
\item statistics collectors;
\item simulation responses;
\item simulation controls;
\item entity groups.
\end{itemize}

This is the core of the language's domain specificity. It is a language for
simulation state, not a generic calculator.

\begin{figure}[htbp]
\centering
% FIGURE-SPEC-BEGIN
% ID: fig:expression-symbol-resolution
% Source of truth:
%   - source/parser/parserBisonFlex/lexerparser.ll
%   - source/parser/parserBisonFlex/bisonparser.yy
%   - source/parser/Genesys++-driver.cpp
% Purpose:
%   Show how a tokenized symbol is resolved to a model object or a runtime value.
% Required visual content:
%   - scanner lookup against the live model
%   - attribute, variable, queue, resource, set, counter, and formula branches
%   - simulation response and control fallback
%   - illegal-token branch
% Accessibility description:
%   A resolution diagram that makes it obvious which names come from the model and which ones become parser errors.
% Validation:
%   The diagram must match the current lookup order and the current illegal-token fallback.
% FIGURE-SPEC-END
\figureplaceholder{figs/generated/expression-symbol-resolution}
\caption{Symbol resolution from scanner token to model-aware value.}
\label{fig:expression-symbol-resolution}
\end{figure}

\section{Indexed access, arrays, and matrices}
\label{sec:expression-reference-indexing}

The language supports sparse indexed access with bracket syntax. The grammar
builds a comma-separated index key and uses it to look up values in the
underlying sparse store.

The practical surface is:
\begin{itemize}
\item \texttt{name};
\item \texttt{name[1]};
\item \texttt{name[1,2]};
\item \texttt{name[1,2,3]};
\item longer sparse keys when the backing data definition supports them.
\end{itemize}

This is how the current code represents array-like and matrix-like data. There
is no separate array literal syntax or matrix constructor in the expression
language. Instead, the language uses indexed model data.

The unit tests cover both reading and writing through this form:
\begin{itemize}
\item variable reads for scalar and indexed values;
\item attribute reads for scalar and indexed values;
\item indexed assignments for variables and attributes;
\item missing sparse entries returning \texttt{0}.
\end{itemize}

\begin{figure}[htbp]
\centering
% FIGURE-SPEC-BEGIN
% ID: fig:expression-context-model
% Source of truth:
%   - source/tests/unit/test_parser_expressions.cpp
%   - source/parser/parserBisonFlex/bisonparser.yy
% Purpose:
%   Show how current-event state, entity ownership, and sparse indices affect expression evaluation.
% Required visual content:
%   - current event and current entity
%   - scalar and indexed attribute lookup
%   - scalar and indexed variable lookup
%   - assignment path for attribute and variable writes
%   - fallback to zero when the indexed slot is absent
% Accessibility description:
%   A context diagram that explains why the same symbol can evaluate differently outside and inside an event.
% Validation:
%   The diagram must reflect the current event-dependent attribute path and the variable store path.
% FIGURE-SPEC-END
\figureplaceholder{figs/generated/expression-context-model}
\caption{Context-dependent evaluation for attributes and variables.}
\label{fig:expression-context-model}
\end{figure}

\section{Errors and recovery}
\label{sec:expression-reference-errors}

The parser distinguishes between illegal characters and unknown literals. It
also surfaces function-specific validation errors for stochastic functions.

The current behavior is:
\begin{itemize}
\item illegal tokens become parse failures;
\item unknown symbols become literal-not-found errors;
\item invalid distribution parameters return recoverable errors when the
      caller uses the non-throwing parse path;
\item the driver stores a last-error message for inspection.
\end{itemize}

The exact message format is intentionally conservative. The important point is
that the parser does not silently accept invalid syntax or invalid stochastic
parameters.

\begin{figure}[htbp]
\centering
% FIGURE-SPEC-BEGIN
% ID: fig:expression-error-anatomy
% Source of truth:
%   - source/parser/Genesys++-driver.cpp
%   - source/parser/parserBisonFlex/bisonparser.yy
%   - source/tests/unit/test_parser_expressions.cpp
% Purpose:
%   Show the current error path from scanner failure or invalid function parameters to the driver error state.
% Required visual content:
%   - illegal character branch
%   - unknown literal branch
%   - invalid probability-parameter branch
%   - thrown versus non-throwing parse path
%   - stored error message
% Accessibility description:
%   A failure-flow diagram that helps the reader distinguish syntax errors from recoverable semantic errors.
% Validation:
%   The diagram must match the current driver error handling and the tested invalid distribution cases.
% FIGURE-SPEC-END
\figureplaceholder{figs/generated/expression-error-anatomy}
\caption{Error anatomy for illegal tokens and invalid parameters.}
\label{fig:expression-error-anatomy}
\end{figure}

\section{Examples reproduced by tests}
\label{sec:expression-reference-examples}

The examples below are all covered by the current unit tests or by direct
evidence in the parser implementation.

\begin{table}[htbp]
\centering
\caption{Current expression examples reproduced by tests.}
\label{tab:expression-reference-examples}
\small
\begin{tabular}{|p{5.2cm}|p{3.0cm}|p{4.7cm}|}
\hline
\textbf{Expression} & \textbf{Expected result} & \textbf{Evidence} \\
\hline
\texttt{1+2*3} & \texttt{7} & precedence test \\
\hline
\texttt{(1+2)*3} & \texttt{9} & precedence test \\
\hline
\verb|2^3^2| & \texttt{512} & right-associative power test \\
\hline
\texttt{1<2 and 3>2} & \texttt{1} & logical precedence test \\
\hline
\texttt{not(1<2)} & \texttt{0} & logical-not test \\
\hline
\texttt{round(3.6)} & \texttt{4} & math-function test \\
\hline
\texttt{ParserVarND[1,2,3,4,5]} & \texttt{14} & sparse variable read test \\
\hline
\texttt{ParserAttrND[1,2,3,4,6]=36} & \texttt{36} & sparse attribute write test \\
\hline
\texttt{unif(0.8,1.0)} & finite value in range & stochastic validity test \\
\hline
\texttt{beta(2,3,1,0)} & failure & invalid-parameter recovery test \\
\hline
\end{tabular}
\end{table}

These examples matter because they show the language as an operational tool.
The reference is not complete unless an expression can be reproduced under the
current parser, driver, and model semantics.

\section{Unsupported syntax and explicit limits}
\label{sec:expression-reference-unsupported}

The current reference does not claim support for the following:
\begin{itemize}
\item quoted string literals;
\item standalone string-expression evaluation via \texttt{val},
      \texttt{eval}, or \texttt{leng};
\item a dedicated array literal syntax;
\item a dedicated matrix literal syntax;
\item recursive formula evaluation through \texttt{FORM[...]};
\item full discrete-distribution support through \texttt{disc(...)};
\item a general-purpose programming language with full \texttt{for} and
      \texttt{else} semantics.
\end{itemize}

Some of those forms are partially tokenized or partially parsed, but the
current implementation does not provide a finished semantic contract for
them. In particular, the grammar still contains clear TODO paths for
\texttt{disc}, formula recursion, and some plugin-backed functions, so those
forms must not be documented as complete features.

\section{Summary}
\label{sec:expression-reference-summary}

The current expression language is a model-aware, sparse-indexed, numeric
language. It supports arithmetic, logic, selected math functions, stochastic
functions, and direct reads and writes against live model data. It also has
clear limits: no quoted strings, no dedicated array or matrix literals, and no
finished recursion for formulas or discrete distributions.

That conservative boundary is the correct reference for the manual today.
